Sections~\ref{ch:1} through~\ref{ch:7} of these Notes are based on lectures I gave at Cornell University in the Spring of 1972. They are expanded from Notes mimeographed at Cornell in May of 1972 and entitled \emph{Automorphic Forms and Representations of Adele Groups}. I am grateful to E.~M.~Stein for suggesting that I expand those Notes for publication by the Princeton University Press and that I incorporate into them the material of Sections~\ref{ch:8} through~\ref{ch:10}. These last three sections are based on lectures I gave at the Institute for Advanced Study, Princeton, in the Spring of 1973. I am indebted to the Institute for its hospitality as well as for the atmosphere it created for serious work.

The subject matter of these Notes is the interplay between the theory of automorphic forms and group representations. One goal is to interpret some recent developments in this area, most significantly the theory of Jacquet-Langlands, working out, whenever possible, explicit consequences and connections with the classical theory. Another goal is to collect as much information as possible concerning the decomposition of $\bL^{2}(\GL(2,\Q)\backslash \GL(2,\A))$. Although each particular section has its own introduction describing the material covered I would like to add the following orienting remarks to this Preface.

Sections~\ref{ch:1} through~\ref{ch:5} are preliminary in nature and their purpose is to spell out the explicit relations between classical cusp forms and certain irreducible constituents of $\bL^{2}(\GL(2,\Q)\backslash \GL(2,\A))$. Here I collect only those facts from representation theory and the classical theory of forms which are crucial to the sequel. Parts of these sections are either new or part of the subject's ``folklore.'' References to the existing literature are to be found in the ``Notes and References'' at the end of each section and individual acknowledgements are made whenever possible.

Sections~\ref{ch:6} through~\ref{ch:10} deal with Jacquet-Langlands' theory and some important questions related to it. Section~\ref{ch:8} describes the continuous spectrum of $\bL^{2}(\GL(2,\Q)\backslash \GL(2,\A))$ and is perhaps the least self-contained. The remaining sections, including Section~\ref{ch:9} on the trace formula, concern the discrete spectrum. I have included a complete proof of the trace formula for $\GL(2)$ primarily because the ideas involved here are still not well known. I also wanted there to be no doubt in the reader's mind that the proof of Jacquet-Langlands' Theorem~\ref{thm:10.5} is now complete. In writing Section~\ref{ch:9} I have followed J.~G.~Arthur's as yet unpublished manuscript on the trace formula for rank one groups and I wish to thank him for allowing me to do so.

Scattered throughout these Notes are some new results and proofs which I have not described elsewhere. I am indebted to my colleagues at Cornell, in particular K.~S.~Brown, W.~H.~J.~Fuchs, A.~W.~Knapp, S.~Lichtenbaum, O.~S.~Rothaus, R.~Stanton, H.~C.~Wang, and W.~C.~Waterhouse, for help and encouragement, and to J.~G.~Arthur, P.~Cartier, W.~Casselman, R.~Howe, R.~Hotta, M.~Karel, R.~P.~Langlands, R.~Parthasarathy, P.~J.~Sally, Jr., and T.~Shintani, for helpful conversations and correspondence related to the results described here. I especially wish to thank R.~P.~Langlands for much valuable information and inspiration.

The first typing of these Notes was done at Cornell by Esther Monroe, Dolores Pendell and Ruth Hymes. Their unusual efficiency and expertise was greatly appreciated.

\begin{flushright}
ITHACA\\[0.3em]
DECEMBER 1973
\end{flushright}
